%
% $Id: $

\documentclass{beamer}
\usetheme{Warsaw}
\usebackgroundtemplate{\includegraphics[width=\paperwidth]{format/libresoft-bg-soft.png}}
\usepackage[spanish]{babel}
\usepackage[latin1]{inputenc}
\usepackage{graphics}
\usepackage{amssymb} % Simbolos matematicos

%\definecolor{libresoftgreen}{RGB}{162,190,43}
%\definecolor{libresoftblue}{RGB}{0,98,143}

%\setbeamercolor{titlelike}{bg=libresoftgreen}

%% Metadatos del PDF.
\hypersetup{
  pdftitle={Background and framework},
  pdfauthor={Jesus M. Gonzalez-Barahona},
  pdfcreator={GSyC/LibreSoft, Universidad Rey Juan Carlos},
  pdfproducer=PDFLaTeX,
  pdfsubject={Master on Libre Software},
}
%%

\begin{document}

\title{Master on Libre Software}
\subtitle{Presentation and Kick-off}
\author{Jesus M. Gonzalez-Barahona, Gregorio Robles}
\institute{\{jgb,grex\}@gsyc.es \\
GSyC/LibreSoft, Universidad Rey Juan Carlos}

\date{Fuenlabrada (Spain), May 23rd, 2012}

\frame{
\maketitle
\begin{center}
\includegraphics[width=6cm]{format/gsyc-urjc}
\end{center}
}

% Si el titulo o el autor se quieren acortar para los pies de p�gina
% se pueden redefinir aqu�:
%\title{Titulo corto}
%\author{Autores abreviado}


% LICENCIA DE REDISTRIBUCION DE LAS TRANSPAS
\frame{
~
\vspace{3cm}

\begin{flushright}
\copyright 2002-2012 Jesus M. Gonzalez-Barahona, Gregorio Robles. \\
Some rights reserved. \\
This presentation is distributed under the \\
``Attribution-ShareAlike 3.0'' license, by Creative Commons, \\
available at \\
\url{http://creativecommons.org/licenses/by-sa/3.0/}

\vspace{.5cm}

The original of this presentation is available in \\
\url{http://master.libresoft.es}
\end{flushright}
}


%\section{}

%%---------------------------------------------------------------

\begin{frame}
\frametitle{Second edition of the official master's degree is starting}

\begin{itemize}
\item Libre software PhD subject\\
  now in a Master's degree on ITC
  since 2003
\item Master's degree on libre software \\
  URJC-specific, ``t�tulo propio'' \\
  thanks to Caixa Nova, Igalia \\
  since 2007
\item Master's degree on libre software, Madrid edition \\
  URJC-specific, ``t�tulo propio'' \\
  2008-2010
\item Official master's degree in Madrid \\
  ``m�ster universitario'' \\
  valid in all the European Higher Education Area \\
  since 2010
\end{itemize}

\end{frame}


%%---------------------------------------------------------------

\begin{frame}
\frametitle{Key design idea(s)}

\begin{center}
{\Large
``To obtain in one year the competences that a current
  expert in libre software acquired by experience during the last 10-15
  years''
}
\end{center}

\begin{itemize}
\item Holistic approach
\item In line with other experiences in Europe and worldwide
\item One step ahead of ``just technical skills'' (but not forgetting
  technical issues)
\item Rooted in real world experiences and academic evidence
\item Let's explore new learning processes
\end{itemize}

\end{frame}

%%---------------------------------------------------------------

\begin{frame}
\frametitle{Some design assumptions}

\begin{itemize}
\item Knowledge about libre software is maturing, facts are available
\item Industry and society need professionals well trained in libre software
\item Target profiles: IT engineers (informatics, telecom, etc.),
  but open to other backgrounds
\item Goal: comprehensive, multidisciplinary education on the
  matter \\
  ``To reach, as a result of the master, the level of an
  expert with good knowledge of the libre software world''
\item All materials available as libre documents
\item Global community: English for written communication
\item Student activities as open as possible
\end{itemize}
\end{frame}


%--------------------------------------------------------
\begin{frame}
\frametitle{Some goals for the master's program}

\begin{itemize}
\item Meeting point around \emph{libre software} in Madrid
\begin{itemize}
  \item Academia
  \item Community
  \item Industry
  \item (Politics)
\end{itemize}
\item The \emph{libre software} way, e.g., synchronous (IRC) and asynchronous (mailing lists, forums) communication channels
\item Service to the \emph{libre software} community with publicly available learning materials
\item Methodology based on blended learning
\end{itemize}

\end{frame}

%--------------------------------------------------------
\begin{frame}
\frametitle{Blended learning}

\begin{itemize}
\item Combination of good practices of on-site and off-site learning
\item Suites well to the geo-chracteristics of the region of Madrid...
\item ...and the lack of time availability of students
\item Theoretical lessons ~25\%: on-site ~80\%
\item Practical lessons ~60\%): on-site ~30\%
\item Evaluation ~15\%: on-site ~30\%
\end{itemize}

\end{frame}

%--------------------------------------------------------
\begin{frame}
\frametitle{Virtual Campus}

\begin{center}
\includegraphics[width=10.5cm]{moodle-master}
\end{center}

Strong use of Moodle to support the learning experience

\begin{flushright}
\url{http://docencia.etsit.urjc.es/moodle/course/view.php?id=135}
\end{flushright}

\end{frame}

%%---------------------------------------------------------------

\begin{frame}
\frametitle{Schedule framework and support activities}

(the final calendar can be downloaded from the web)

\begin{itemize}
\item On-site sessions: 
  \begin{itemize}
  \item Thursday, 17:00-21:00
  \item Friday, 17:00-21:00
  \end{itemize}
\item On-site short seminars:
  \begin{itemize}
  \item Thursday, 16:00-17:00
  \item Video recorded, schedule will be published
  \end{itemize}
\item On-site mentoring sessions on demand
\item On-line mentoring (IRC/labs):
  \begin{itemize}
  \item Wednesday, 19:00-21:00 (on demand)
  \item irc.freenode.net \#mswl
  \end{itemize}
\item Continuous support on Moodle, usual in IRC
\item Syllabus document for each subject, with scheduling details
\end{itemize}
\end{frame}


%%---------------------------------------------------------------

\begin{frame}
\frametitle{Cross-section activities}

\begin{itemize}
\item Learning blogs (limited only by topic)
\item Collaborative notebook
\item Identi.ca / twitter channel \#mswl
\item External seminars
\item ...
\end{itemize}
\end{frame}

%--------------------------------------------------------
\begin{frame}
\frametitle{Structure of the program}

General structure (September - May):

\begin{itemize}
\item 30 ECTS of mandatory teaching subjects
\item 6 ECTS of elective teaching subjects
\item 24 ECTS of Practicum (12) and Final Project (12)
\item Most subjects: 3 ECTS
\end{itemize}

By semester:

\begin{itemize}
\item First: 18 ECTS mandatory
\item Second: 12 ECTS mandatory, 6 ECTS elective
\item Practicum and Project spans across both semesters
\end{itemize}
\end{frame}

%%---------------------------------------------------------------

\begin{frame}
\frametitle{First semester}

Mandatory:

\begin{itemize}
\item Introduction: 3 credits
\item Legal aspects: 3 credits
\item Economic aspects: 3 credits
\item Developers and motivation: 3 credits
\item Development and tools: 3 credits
\item Project evaluation: 3 credits
\end{itemize}
\end{frame}

%%---------------------------------------------------------------

\begin{frame}
\frametitle{Second semester}

Mandatory:

\begin{itemize}
\item Case studies I: 6 credits
\item Project management: 3 credits
\item Case studies II: 3 credits
\end{itemize}

Elective (6 credits):

\begin{itemize}
\item Development and sysadmin track
  \begin{itemize}
  \item Advanced development: 3 credits
  \item System integration: 3 credits
  \end{itemize}
\item Applications and communities
  \begin{itemize}
  \item Deployment: 3 credits
  \item Communities: 3 credits
  \end{itemize}
\end{itemize}

\end{frame}



%%---------------------------------------------------------------

\begin{frame}
\frametitle{Subject: Introduction}

\begin{itemize}
\item Lecturers:
  \begin{itemize}
  \item Gregorio Robles (coord.)
  \item Jesus M. Gonzalez-Barahona
  \end{itemize}
\item General concepts about libre (free, open source) software
\item Pre-requisite for almost all the subjects
\item Quite general issues: details in other subjects
\item Should provide a general framework
\item Prone to discussion
\item Base text: \\
  ``Introduction to free software'' \\
  Jesus M. Gonzalez-Barahona, Joaquin Seoane, Gregorio Robles
\end{itemize}
\end{frame}

%%---------------------------------------------------------------

\begin{frame}
\frametitle{Subject: Legal Aspects}

\begin{itemize}
\item Lecturers:
   \begin{itemize}
   \item Miguel Vidal (coord.)
   \item Gregorio Robles
   \end{itemize}
\item Why?: Licenses are legal basis and ``social contract'' of
libre software.
\item Main topics:
   \begin{itemize}
    \item Intellectual property (IP) law for non-lawyers: Copyright,
Patents, Trademarks, Warranties...
    \item How FLOSS and IP work together
    \item Free and Open Source Software licenses (permissive and reciprocal)
    \item Free cultural works (documentation, images, Creative Commons
licenses)
    \item Case studies and practical issues (choosing a license, 
compatibility,
forking, merging/mixing... myths)
   \end{itemize}
\end{itemize}

\end{frame}

%%---------------------------------------------------------------

\begin{frame}
\frametitle{Subject: Economic Issues}

\begin{itemize}  
\item Lecturers:
  \begin{itemize}
  \item Jesus M. Gonzalez-Barahona (coord.)
  \item David Lopez
  \end{itemize}
\item Main topics:
  \begin{itemize}
  \item Economy of libre software
  \item Related business and business models
  \item Impact of business on licensing, development (and viceversa)
  \item For-profit and volunteers collaboration
  \item Impact of libre software in companies and strategies
  \end{itemize}
\end{itemize}

\end{frame}

\begin{frame}
\frametitle{Subject: Developers and motivation}

\begin{itemize}
  \item Lecturers:
   \begin{itemize}
   \item Felipe Ortega (coord.)
   \item Miguel Vidal
   \item Juan Jos� Amor
\end{itemize}
\item Main topics:
  \begin{itemize}
   \item Motivations of developers in libre software projects.
   \item Roles assumed by developers in libre software projects.
   \item Leadership and sustainability of libre software projects.
   \item Social structure of libre software communities.
  \end{itemize}

\end{itemize}

\end{frame}


%%---------------------------------------------------------------
%%---------------------------------------------------------------
\begin{frame}
\frametitle{Subject: Development and tools}

\begin{itemize}
  \item Lecturers:
   \begin{itemize}
   \item Israel Herr\'aiz (coord.)
   \item Micael Gallego
   \item Francisco Gort�zar
   \item Carlos Garc�a Campos
\end{itemize}
\item Main topics:
  \begin{itemize}
   \item Tools and development environment
     \begin{itemize}
     \item The version control system: Git
     \item Integrated Development Environments (Eclipse)
     \end{itemize}
   \item Python programming language
   \item Building desktop application with the GNOME development
     platform
  \end{itemize}
\end{itemize}

\end{frame}

%%---------------------------------------------------------------
\begin{frame}
\frametitle{Subject: Project evaluation}

\begin{itemize}
\item Lecturers:
  \begin{itemize}
  \item Daniel Izquierdo (coord.)
  \item Pedro Coca
  \item Pedro Garc�a
  \end{itemize}
\item Main topics:
  \begin{itemize}
  \item Introduction to data sources and projects' evaluation.
  \item Light-weight (OpenBRR) and heavy-weight (QualOSS) evaluation 
    models .
  \item Meta repositories of information
  \item Quality assurance: Mozilla community case study
  \end{itemize}
\end{itemize}

\end{frame}

%%---------------------------------------------------------------
\begin{frame}
\frametitle{Subject: Case studies I, II}

\begin{itemize}
   \item Lecturers (coordination):
    \begin{itemize}
    \item Pedro Garc�a (coord.)
\end{itemize}
\item Case studies I Seminars (Thursday, 16:00-17:00), 
  first term, recorded in video.
\item Case studies II Seminars (Friday, 17:00-21:00),
  second term, recorded in video.
\item Other on-line activities
\item Main topics:
   \begin{itemize}
   \item Cross-section issues for all the master
   \item Complementary talks
   \item First-hand experiences and case studies
   \end{itemize}
\item Students required to document the topics covered
\item Essays on specific topics related to the seminars
\end{itemize}

\end{frame}

%%---------------------------------------------------------------
\begin{frame}
\frametitle{Subject: Project management}

\begin{itemize}
\item Lecturers:
  \begin{itemize}
  \item Felipe Ortega (coord.)
  \item Santiago Gala
  \item Teo Romera
  \end{itemize}
\item Main topics:
  \begin{itemize}
  \item Factors influencing management of libre software projects.
  \item How to set up your project and nurture participation.
  \item How to integrate yourself in an existing libre software project.
  \item Netiquette and interaction recommendations.
  \item Community management and volunteers.
  \end{itemize}
\end{itemize}

\end{frame}

%-------------------------------------------------------------------
\begin{frame}
\frametitle{Subject: Systems Integration (elective)}

\begin{itemize}
   \item Lecturers:
     \begin{itemize}
       \item Jose Castro (coord.)
       \item Miguel Vidal
     \end{itemize}
\item State-of-the-art system technologies overview
\item Tasks that sysadmins are expected to perform with libre software
\item Main topics:
   \begin{itemize}
   \item Essential tasks of systems administration. Skills, policies and
procedures.
    \item Server-side Internet services (heterogeneous Unix-like systems 
deployment)
    \item Virtualization technologies, clusters HA and cloud computing 
(IaaS).
   \end{itemize}
\end{itemize}

\end{frame}

%-------------------------------------------------------------------
\begin{frame}
\frametitle{Subject: Deployment (elective)}

\begin{itemize}  
  \item Lecturers:
    \begin{itemize}
      \item Jos� Gato (coord.)
      \item Miguel Vidal
    \end{itemize}
\item Migration to libre software environments
\item Deployment of libre software with a minimal negative impact
\item Scheduling, scoping and goals of deployment projects
\item Requirements of software and users
\item Analysis of different options: review of different application domains
\item Training and Total Cost of Ownership (TCO)
\end{itemize}

\end{frame}

%-------------------------------------------------------------------
\begin{frame}
  \frametitle{Subject: Communities (elective)}

  \begin{itemize}
  \item Lecturers
    \begin{itemize}
    \item Daniel Izquierdo (coord.)
    \item Felipe Ortega
    \end{itemize}
  \item Main topics:
    \begin{itemize}
    \item Introduction to FLOSS communities and data mining approach
    \item Cases of study:
      \begin{itemize}
      \item Measuring activity in FLOSS projects and effort distribution
      \item Developer turnover and knowledge management
      \item Territoriality and experience
      \end{itemize}
    \end{itemize}
  \end{itemize}
\end{frame}

%-------------------------------------------------------------------
\begin{frame}
\frametitle{Subject: Advanced Development (elective)}

\begin{itemize}  
  \item Lecturers:
    \begin{itemize}
      \item Roberto Calvo (coord.)
    \end{itemize}

\item Android development
\item History \& philosophy of Android project
\item Eclipse \& SDK
\item Application Environment: Layouts, Activity
\item Maps \& Geolocation
\item Upload Apps to Market and Advertising policies
\item Appendix: Augmented reality \& Arduino

\end{itemize}

\end{frame}

%-------------------------------------------------------------------
\begin{frame}
\frametitle{Subject: Practicum}

\begin{itemize}
\item Practicum on Libre Software
\item Coordinator: Gregorio Robles
\item Mentoring meetings instead of lectures
\item Aimed to provide a hands-on / real-world experience in libre
  software, either in companies or in projects
\item Can be related to the final project, but that is not mandatory
\item Could be linked to professional activities if related to libre
software
\item Decide and schedule asap!
\end{itemize}

\end{frame}

%%---------------------------------------------------------------

\begin{frame}
\frametitle{Subject: Final project}

\begin{itemize}  
\item Lecturer (coordination):
  \begin{itemize}
  \item Gregorio Robles
  \end{itemize}
\item Two chances:
  \begin{itemize}
  \item General proposal: in-depth analysis of libre software scenarios
  \item Customized proposals: suited to specific needs/wills
  \end{itemize}
\item Results will be public (as most other materials for the degree)
%\item Video presentation with main results
\item Important: start with it asap
\end{itemize}

\end{frame}

%-------------------------------------------------------------------
\begin{frame}
\frametitle{The labs}

\begin{itemize}
\item Laboratorios II building, ``Linux labs''
\item Up-to-date Ubuntu
\item Networked environment, common homes, need account
\item Can also bring your own laptops (recommended for some subjects)
\item Could install the same environment in laptops or home computers
\item Some classes will be there
\end{itemize}

\begin{flushright}
\url{http://bilo.gsyc.es}
\end{flushright}

\end{frame}

%-------------------------------------------------------------------
\begin{frame}
\frametitle{Computers in class}

\begin{itemize}
\item Please, bring your laptops with you...
\item ...and use them while in class:
  \begin{itemize}
  \item Write for the notebook
  \item Browse references in the web
  \item Microblog
  \item Download documents, write links for later use
  \end{itemize}
\item For practice work, either your laptop or the labs
\item Ubuntu is the preferred platform
\end{itemize}

\vspace{.5cm}

\begin{center}
Remember, you can install (and use)  everything we talk about, \\
this is libre software
\end{center}

\end{frame}

%%---------------------------------------------------------------

%% \begin{frame}
%% \frametitle{Experience with non-enrolled students}

%% \begin{itemize}
%% \item Translating practices from the libre software community to education
%% \item Peer-to-peer learning, assisted by lecturers
%% \item Some commitments, no title (nor degree) at the end
%% \item Three subjects: introduction, legal aspects, economic aspects
%% \item Temptative experience, could be the beginning of a way to explore
%% \end{itemize}
%% \end{frame}

%%---------------------------------------------------------------

\begin{frame}
\frametitle{Fuenlabrada campus}

\begin{center}
\includegraphics[width=10.5cm]{campus-fuenla}
\end{center}

\end{frame}


%%---------------------------------------------------------------

\begin{frame}

\begin{center}
{\Huge Enjoy!}
\end{center}

\vspace{3cm}

\begin{flushright}
\url{http://master.libresoft.es/}
\end{flushright}

\end{frame}






\end{document}
